\section{Cyberbezpieczeństwo i ochrona danych}
\subsection*{Dostęp do systemów}
Każdy operator ma imienne konto. Role rozdzielają podgląd, eksport, nadawanie kart i administrację. Konta serwisowe są aktywowane na ograniczony czas, a dostęp zdalny wymaga VPN i MFA. Kamery i centrale nie są publikowane w Internecie; wyłącza się nieużywane usługi i funkcje chmurowe producenta. Połączenia szyfruje się tam, gdzie wspiera to dobrany system.

Ruch między VLAN jest domyślnie blokowany. Lista wyjątków obejmuje konkretne źródła, cele i usługi. Urządzenia korzystają ze wspólnego lokalnego źródła czasu. Logi są kopiowane do wydzielonej usługi z kontrolą dostępu; retencja 90 dni jest założeniem roboczym wymagającym uzasadnienia, a nie wymogiem ustawowym [9].

\subsection*{Dane osobowe i retencja}
Obraz może identyfikować osoby, więc administrator musi określić cel, podstawę prawną, niezbędny obszar obserwacji i okres przechowywania. Przyjęte do bilansu \textbf{7 dni} nie jest wymaganiem prawa. Należy uzasadnić ten okres i skrócić go, jeśli cel można osiągnąć krócej. Wyodrębniony materiał incydentu ma oddzielną, udokumentowaną podstawę i termin usunięcia [5].

Przed wdrożeniem przewidziano ocenę skutków dla ochrony danych dla monitoringu dużej strefy publicznej, konsultację IOD oraz uwzględnienie właściwych przepisów sektorowych i prawa pracy. W koncepcji przyjęto oznaczenie stref monitorowanych, pełną informację dla osób, maskowanie obszarów zbędnych i brak rejestracji dźwięku. Nie projektuje się kamer w toaletach ani pomieszczeniach wymagających szczególnej prywatności [5].

\subsection*{Kopie i aktualizacje}
Konfiguracje central, przełączników, VMS i KD kopiuje się po każdej zatwierdzonej zmianie. Raz na kwartał wykonuje się próbę odtworzenia kopii w odizolowanym środowisku. Aktualizacje najpierw przechodzą test zgodności, a następnie są wdrażane w uzgodnionym oknie z możliwością powrotu do poprzedniej wersji. Dla SSP zakres zmian określa producent i uprawniony serwis.

\begin{note}{Co trafia do publicznego repozytorium}
Wyłącznie fikcyjna geometria, przykładowe identyfikatory, treść akademicka i źródła rysunków. Rzeczywistych konfiguracji, haseł, kluczy, nagrań i planów ochrony nie należy dodawać przy późniejszej adaptacji projektu do realnego obiektu.
\end{note}
\src{[5] EROD -- monitoring i dane osobowe; [9] CISA -- segmentacja i ograniczanie dostępu}
